\documentclass[11pt]{article}
\usepackage{fontspec}
\setmainfont{Times New Roman}
\setsansfont{Arial}
\setmonofont{Consolas}
\usepackage{amsmath,amssymb}
\usepackage{geometry}
\usepackage{microtype}
\geometry{a4paper,margin=3cm}
\setlength{\parindent}{1.25em}
\setlength{\parskip}{0pt}
\emergencystretch=4em
\allowdisplaybreaks

\begin{document}

\begin{center}
\textbf{\S\ 8. \quad Formy tretjego poręda (Sistem III).}
\end{center}

\noindent\textbf{Svijanje $f$ s $\nu$.}

\noindent\textit{Nereducibilne formy:}
\[
\begin{array}{cccc}
a_\nu&\cdot&\cdot&\cdot\\
\cdot&a_\nu^2&\cdot&\cdot\\
\cdot&\cdot&\cdot&\cdot\\
\cdot&\cdot&\cdot&a_\nu^4=i.
\end{array}
\]

\noindent\textit{Redukcija:}
\[
a_\nu^3=\frac13 i u_x\qquad\text{(formula (2.)).}
\]
\noindent\textit{Poslědki:} 1) $a_\nu^3$ jest reducent. 2) $f$ vhodi v svijanje s $\nu$ samo v produktah s kontragredientnymi formami $(j)$. To samo važi za $\nu$ v odnosu k $f$.

\end{document}
